\clearpage
\VerdierPhysicalPageAnchor{0225}%
\begingroup\small\setlength{\parskip}{0pt}\setlength{\emergencystretch}{1.5em}
\setlength{\abovedisplayskip}{4pt plus 2pt minus 2pt}%
\setlength{\belowdisplayskip}{4pt plus 2pt minus 2pt}%

\VerdierSemanticAnchor{ch4d-4-1-1-cont}
こうして関手
$\mathcal{F}\ell_1(\widetilde{\mathbb Z})\longrightarrow K(\mathcal A)$
が定義された。

$(p,q,r)$ を $\mathcal{F}\ell_2(\widetilde{\mathbb Z})$ の元とする。

\VerdierSemanticAnchor{ch4d-4-1-1-6}
\begin{equation}
\delta(p,q,r)\colon\mathcal Y(q,r)\longrightarrow\mathcal Y(p,q)[1],
\tag{4.1.1.6}
\end{equation}
を、次の複体の射のホモトピー類として定義する：

\VerdierSemanticAnchor{ch4d-4-1-1-7}
\begin{equation}
d(p,q,r)\colon\mathcal Y(q,r)\longrightarrow\mathcal Y(p,q)[1],
\tag{4.1.1.7}
\end{equation}
この射は次の式で定義される：

\VerdierSemanticAnchor{ch4d-4-1-1-8}
\begin{equation}
d(p,q,r)^i=
\begin{pmatrix}
0 & \operatorname{id}_{Y(q)}^{i+1}\\
0 & 0
\end{pmatrix}.
\tag{4.1.1.8}
\end{equation}

（第 I 章、3.3.6）により、$K(\mathcal A)$ の次の三角

\VerdierSemanticAnchor{ch4d-4-1-1-distinguished-triangle}
\begin{center}
\resizebox{\textwidth}{!}{%
$\bigl(\mathcal Y(p,q),\mathcal Y(p,r),\mathcal Y(q,r),
\Phi((p,q),(p,r)),\Phi((p,r),(q,r)),\delta(p,q,r)\bigr)$%
}
\end{center}
は完全三角であり、$(p,q,r)$ に関手的に依存することが知られている。
したがって $K(\mathcal A)$ に値をとるスペクトル対象
$(\mathcal Y(p,q),\delta(p,q,r))$ が構成された。

$Y\colon\widetilde{\mathbb Z}\longrightarrow
\operatorname{comp}(\mathcal A)$ と
$Y'\colon\widetilde{\mathbb Z}\longrightarrow
\operatorname{comp}(\mathcal A)$ を二つの関手とし、
$m\colon Y\longrightarrow Y'$ を関手の射とする。射 $m$ は付随する
スペクトル対象の射を誘導する：

\VerdierSemanticAnchor{ch4d-4-1-1-9}
\begin{equation}
\mathcal M(p,q)\colon\mathcal Y(p,q)\longrightarrow\mathcal Y'(p,q).
\tag{4.1.1.9}
\end{equation}

\VerdierSemanticAnchor{ch4d-4-1-2}
\medskip\noindent\textbf{4.1.2.}\quad
$Y$ を $\operatorname{comp}(\mathcal A)$ の対象とする。次のように置く：

\VerdierSemanticAnchor{ch4d-4-1-2-1}
\begin{equation}
q\in\mathbb Z,\qquad
Y(q)=0\longrightarrow0\longrightarrow Y^{-q}\longrightarrow
Y^{-q+1}\longrightarrow Y^{-q+2}\longrightarrow\cdots,
\tag{4.1.2.1}
\end{equation}
ここで、この複体の微分は $Y$ の微分から誘導される。同様に、

\VerdierSemanticAnchor{ch4d-4-1-2-2}
\begin{equation}
Y(-\infty)=0,\qquad Y(+\infty)=Y.
\tag{4.1.2.2}
\end{equation}
と置く。モノ射

\VerdierSemanticAnchor{ch4d-4-1-2-filtration-chain}
\[
0=Y(-\infty)\hookrightarrow\cdots\hookrightarrow Y(q)
\hookrightarrow Y(q+1)\hookrightarrow Y(q+2)\hookrightarrow\cdots
\hookrightarrow Y(+\infty)=Y,
\]
がある。
\endgroup
